\documentclass[12pt,a4paper]{article}
\usepackage[utf8]{inputenc}
\usepackage[T1]{fontenc}
\usepackage[bahasa]{babel}
\usepackage[top=2.5cm,bottom=2.5cm,left=2.5cm,right=2.5cm,headheight=14pt]{geometry}
\usepackage{lmodern}
\usepackage{graphicx}
\usepackage{hyperref}
\usepackage{booktabs}
\usepackage{longtable}
\usepackage{array}
\usepackage{xcolor}
\usepackage{tikz}
\usetikzlibrary{shapes.geometric,arrows.meta,positioning,calc,fit,backgrounds}
\usepackage{pgfplots}
\pgfplotsset{compat=1.18}
\usepackage{listings}
\usepackage{enumitem}
\usepackage{fancyhdr}
\usepackage{titlesec}
\usepackage{setspace}
\usepackage{microtype}
\usepackage{amsmath}
\usepackage{float}
\usepackage{caption}
\usepackage{subcaption}
\usepackage{multirow}
\usepackage{tabularx}
\usepackage{colortbl}

\definecolor{verfinblue}{RGB}{30,64,175}
\definecolor{verifindark}{RGB}{15,23,42}
\definecolor{verfinred}{RGB}{220,38,38}
\definecolor{verfinyellow}{RGB}{202,138,4}
\definecolor{verfingreen}{RGB}{22,163,74}
\definecolor{verfingray}{RGB}{248,250,252}
\definecolor{codebg}{RGB}{245,245,245}

\hypersetup{colorlinks=true,linkcolor=verfinblue,urlcolor=verfinblue,citecolor=verfinblue,pdfauthor={Tim Verifin UGM},pdftitle={Proposal GEMASTIK XIX -- Verifin}}

\pagestyle{fancy}
\fancyhf{}
\fancyhead[L]{\small\textbf{Verifin}}
\fancyhead[C]{\small GEMASTIK XIX 2026}
\fancyhead[R]{\small Universitas Gadjah Mada}
\fancyfoot[C]{\thepage}
\renewcommand{\headrulewidth}{0.4pt}

\titleformat{\section}{\Large\bfseries\color{verfinblue}}{}{0em}{}[\titlerule]
\titleformat{\subsection}{\large\bfseries\color{verifindark}}{}{0em}{}
\titleformat{\subsubsection}{\normalsize\bfseries}{}{0em}{}

\lstset{basicstyle=\small\ttfamily,breaklines=true,backgroundcolor=\color{codebg},frame=single,framesep=4pt,rulecolor=\color{gray!40},xleftmargin=10pt,xrightmargin=10pt,keywordstyle=\bfseries\color{verfinblue},commentstyle=\color{gray},stringstyle=\color{verfingreen},numbers=left,numberstyle=\tiny\color{gray},numbersep=6pt,showstringspaces=false}

\onehalfspacing

\begin{document}

\begin{titlepage}
\thispagestyle{empty}
\begin{center}
\vspace*{0.5cm}
{\color{verfinblue}\rule{\linewidth}{3pt}}
\vspace{0.4cm}
{\small\bfseries\color{verifindark} PAGELARAN MAHASISWA NASIONAL BIDANG TEKNOLOGI INFORMASI DAN KOMUNIKASI}\\[0.25cm]
{\Huge\bfseries\color{verfinblue} GEMASTIK XIX}\\[0.3cm]
{\large\bfseries\color{verifindark} DIVISI PENGEMBANGAN PERANGKAT LUNAK}
\vspace{0.4cm}
{\color{verfinblue}\rule{\linewidth}{1pt}}
\vspace{2cm}
{\fontsize{40}{48}\selectfont\bfseries\color{verfinblue} Verifin}\\[0.4cm]
{\large\itshape\color{verifindark} Verifikasi Lowongan Kerja Berbasis AI, OSINT, dan Explainable AI}
\vspace{2cm}
\begin{tabular}{ll}
\toprule
\multicolumn{2}{c}{\bfseries\large Anggota Tim}\\
\midrule
Ketua Tim & Hafidz Rizqullah Prasetya (24/535493/SV/24243)\\
Anggota 1 & Matthew Hayunaji Priantara (24/536179/SV/24400)\\
Anggota 2 & Akmal Manggala Putra (24/536182/SV/24402)\\
\bottomrule
\end{tabular}
\vspace{1.5cm}
\begin{tabular}{ll}
Sekolah & Sekolah Vokasi\\
Program Studi & Teknologi Rekayasa Perangkat Lunak (D4)\\
Universitas & Universitas Gadjah Mada\\
Kota & Yogyakarta\\
Tahun & 2026\\
\end{tabular}
\vspace{2cm}
{\color{verfinblue}\rule{\linewidth}{3pt}}
\end{center}
\end{titlepage}

\newpage
\tableofcontents
\newpage

\section*{A.\quad Latar Belakang Ide Perangkat Lunak}
\addcontentsline{toc}{section}{A. Latar Belakang Ide Perangkat Lunak}

Tingkat pengangguran terbuka di Indonesia masih menjadi salah satu tantangan ekonomi nasional yang
signifikan. Berdasarkan data Badan Pusat Statistik (BPS), per Februari 2025 jumlah pengangguran
terbuka tercatat mencapai 7,28 juta orang, meningkat dibandingkan periode sebelumnya [1]. Di tengah
tekanan tersebut, jutaan pencari kerja Indonesia setiap harinya menerima tawaran lowongan dari
berbagai kanal yang tidak terverifikasi: Instagram, WhatsApp, Telegram, grup Facebook, hingga papan
pengumuman fisik. Platform-platform ini menyebarkan lowongan secara masif, namun \textbf{tidak ada
satu pun mekanisme yang menjawab pertanyaan mendasar: ``Apakah lowongan ini layak dipercaya?''}

Akibatnya, pencari kerja harus melakukan investigasi manual secara mandiri: mencari nama perusahaan
di Google, mengecek nomor HP di direktori, melihat lokasi di Maps, menelusuri ulasan di media
sosial. Proses ini melelahkan, tidak sistematis, dan sering kali terlambat karena korban baru
menyadari penipuan setelah menyerahkan data pribadi atau uang. Menurut laporan \textit{Global State
of Scams} oleh Global Anti-Scam Alliance (GASA), Indonesia menempati salah satu peringkat kerentanan
tertinggi di Asia Tenggara dengan \textbf{49\% penduduk mengalami kerugian akibat penipuan} [2].
Lebih jauh, krisis ini tidak hanya berujung pada kerugian finansial, melainkan telah bergeser menjadi
pintu masuk Tindak Pidana Perdagangan Orang (TPPO): ribuan pencari kerja terjebak iklan pekerjaan
palsu di luar negeri dan dipaksa bekerja sebagai operator penipuan siber (\textit{scam center}) [3].

Celah yang ada saat ini bukan sekadar permasalahan teknologi---melainkan kegagalan ekosistem. Platform
media sosial tidak bertanggung jawab atas konten lowongan yang disebarkan di dalamnya. Job board
berbayar memiliki proses verifikasi terbatas. Adapun platform gratis sama sekali tidak memiliki
mekanisme kuratif. Akibatnya, \textbf{beban pembuktian jatuh sepenuhnya kepada individu yang paling
rentan}: pencari kerja yang sedang dalam kondisi terdesak secara finansial.

Solusi yang sudah ada---seperti fitur pelaporan di platform media sosial atau imbauan dari Kementerian
Ketenagakerjaan---bersifat reaktif, lambat, dan bergantung pada pelaporan manual. Tidak ada sistem yang
mampu menganalisis sebuah teks lowongan secara real-time dan memberikan penilaian berbasis bukti kepada
pengguna \textit{sebelum} mereka memutuskan untuk merespons.

\subsection*{Inilah celah yang Verifin hadir untuk menutupnya.}

\textbf{Verifin} (\textit{Verifikasi Lowongan Kerja}) adalah platform intelijen berbasis web yang
menganalisis teks atau URL lowongan kerja secara otomatis menggunakan rantai teknologi berlapis:
Named Entity Recognition (NER), klasifikasi NLP, investigasi OSINT multi-sumber, sintesis naratif LLM,
dan penjelasan kausal berbasis Explainable AI (XAI). Keluaran sistem adalah sebuah \textit{Trust Score}
0--100 beserta verdict tiga tingkat (AMAN/WASPADA/BAHAYA) yang didukung oleh penjelasan berbasis bukti
yang dapat diverifikasi oleh pengguna awam.

\vspace{0.5cm}
\textbf{Relevansi dengan GEMASTIK XIX --- Divisi Pengembangan Perangkat Lunak:}

Verifin secara langsung menjawab tantangan nasional yang nyata dengan memanfaatkan kombinasi
mutakhir teknologi AI yang relevan: NLP berbahasa Indonesia, OSINT otomatis, LLM API, dan XAI
untuk transparansi. Sistem ini bukan sekadar demonstrasi teknologi, melainkan perangkat lunak
fungsional dengan nilai guna langsung bagi jutaan pencari kerja Indonesia.

\vspace{0.5cm}

% Infografis: Statistik Kerentanan
\begin{figure}[H]
\centering
\begin{tikzpicture}
\begin{axis}[
  ybar,
  width=13cm, height=6cm,
  bar width=1.2cm,
  xlabel={Kanal Penyebaran Lowongan Palsu},
  ylabel={Persentase (\%)},
  symbolic x coords={WhatsApp,Instagram,Telegram,Facebook,Email,Lainnya},
  xtick=data,
  ymin=0, ymax=70,
  nodes near coords,
  nodes near coords align={vertical},
  every node near coord/.append style={font=\small},
  title={\textbf{Distribusi Kanal Penyebaran Lowongan Penipuan di Indonesia}},
  title style={font=\bfseries\small},
  tick label style={font=\small},
  label style={font=\small},
  grid=major, grid style={dashed,gray!30},
  axis background/.style={fill=verfingray!30},
]
\addplot[fill=verfinblue!80,draw=verfinblue] coordinates {
  (WhatsApp,58) (Instagram,52) (Telegram,38) (Facebook,34) (Email,18) (Lainnya,12)
};
\end{axis}
\end{tikzpicture}
\caption{Distribusi kanal penyebaran lowongan penipuan berdasarkan survei GS Lab \& GASA 2024 [2]}
\end{figure}

\begin{figure}[H]
\centering
\begin{tikzpicture}
\begin{axis}[
  xbar,
  width=12cm, height=4cm,
  bar width=0.7cm,
  xlabel={Persentase Pencari Kerja (\%)},
  symbolic y coords={Terekspos Penipuan Lowongan},
  ytick=data,
  xmin=0, xmax=70,
  nodes near coords,
  nodes near coords align={horizontal},
  title={\textbf{49\% Penduduk Indonesia Pernah Terekspos Penipuan (GASA 2024)}},
  title style={font=\bfseries\small},
  tick label style={font=\small},
  grid=major, grid style={dashed,gray!30},
]
\addplot[fill=verfinred!80,draw=verfinred] coordinates {(49,Terekspos Penipuan Lowongan)};
\end{axis}
\end{tikzpicture}
\caption{49\% penduduk Indonesia mengalami kerugian akibat penipuan (Global Anti-Scam Alliance, 2024) [2]}
\end{figure}


\section*{B.\quad Tujuan dan Manfaat Dikembangkannya Perangkat Lunak}
\addcontentsline{toc}{section}{B. Tujuan dan Manfaat Dikembangkannya Perangkat Lunak}

\subsection*{B.1 Tujuan}

\begin{enumerate}[leftmargin=*, label=\textbf{\arabic*.}]
\item Membangun platform web \textbf{Verifin} yang mampu menganalisis teks atau URL lowongan kerja
      secara otomatis dan menghasilkan \textit{Trust Score} 0--100 beserta verdict AMAN/WASPADA/BAHAYA
      dalam waktu di bawah 30 detik.

\item Mengimplementasikan pipeline analisis berlapis yang mengintegrasikan Named Entity Recognition
      (NER) berbahasa Indonesia, klasifikasi NLP berbasis IndoBERT, investigasi OSINT multi-sumber
      (domain, nomor telepon, perusahaan), sintesis naratif berbasis LLM (GPT-4o-mini via
      OpenAgentic), dan penjelasan kausal berbasis custom SHAP-inspired XAI.

\item Membangun \textit{Fraud Network Graph} berbasis database PostgreSQL yang dapat mendeteksi
      penggunaan ulang entitas-entitas mencurigakan (nomor telepon, domain, nama perusahaan) lintas
      lowongan.

\item Menyediakan fitur \textit{community monitoring} yang memungkinkan pengguna melaporkan dan
      memberikan penilaian terhadap lowongan mencurigakan secara kolektif.

\item Menyediakan REST API publik terdokumentasi untuk memungkinkan integrasi dengan platform
      pihak ketiga.
\end{enumerate}

\subsection*{B.2 Manfaat}

\begin{enumerate}[leftmargin=*, label=\textbf{\arabic*.}]
\item \textbf{Bagi Pencari Kerja:} Mendapat alat verifikasi gratis, cepat, dan berbasis bukti yang
      mengurangi risiko menjadi korban penipuan lowongan. Tidak perlu lagi melakukan investigasi
      manual yang melelahkan.

\item \textbf{Bagi Ekosistem Ketenagakerjaan:} Menciptakan lapisan keamanan kolektif melalui
      \textit{Fraud Network Graph} yang terus berkembang. Setiap analisis baru memperkaya database
      entitas mencurigakan untuk melindungi pengguna berikutnya.

\item \textbf{Bagi Masyarakat Luas:} Berkontribusi dalam penanggulangan TPPO berbasis
      \textit{cyber-recruitment} dengan memutus rantai penipuan di tahap paling awal: sebelum
      korban merespons iklan.

\item \textbf{Bagi Pengembang dan Peneliti:} API publik memungkinkan integrasi ke dalam platform
      lain (ekstensi browser, bot Telegram, job board). Data agregat dari sistem dapat menjadi
      sumber penelitian tentang pola penipuan kerja di Indonesia.

\item \textbf{Bagi Kemajuan NLP Bahasa Indonesia:} Pengembangan sistem NER dan klasifikasi NLP
      yang dioptimalkan untuk domain lowongan kerja berbahasa Indonesia berkontribusi pada
      ekosistem NLP Bahasa Indonesia yang masih terus berkembang [9].
\end{enumerate}


\section*{C.\quad Batasan Perangkat Lunak yang Dikembangkan}
\addcontentsline{toc}{section}{C. Batasan Perangkat Lunak yang Dikembangkan}

Untuk memastikan sistem dapat dibangun, diuji, dan didemonstrasikan secara nyata dalam kerangka
kompetisi, Verifin menetapkan batasan-batasan berikut:

\begin{enumerate}[leftmargin=*, label=\textbf{C\arabic*.}]

\item \textbf{Bahasa Input:}
Sistem dioptimalkan untuk memproses teks lowongan dalam Bahasa Indonesia. Input dalam bahasa
lain (Inggris, campuran) dapat diterima namun akurasi NER dan klasifikasi NLP tidak dijamin
setara.

\item \textbf{Sumber OSINT:}
Investigasi OSINT terbatas pada sumber yang dapat diakses secara publik dan legal:
\begin{itemize}
  \item Whois/RDAP untuk domain
  \item Validasi format dan prefix nomor telepon Indonesia
  \item Pencarian nama perusahaan di database AHU (Administrasi Hukum Umum) via scraping terbatas
  \item Pencarian Google Custom Search API dengan kueri terstruktur
  \item Pengecekan reputasi domain via VirusTotal API (tier gratis)
\end{itemize}
Sistem \textbf{tidak} mengakses database kepolisian, Dukcapil, atau sistem pemerintah lainnya
yang memerlukan otorisasi khusus.

\item \textbf{Skala Database Awal:}
\textit{Fraud Network Graph} pada versi demonstrasi akan diinisialisasi dengan data sintetis
yang representatif. Database komunitas nyata akan terbangun secara organik setelah sistem
diluncurkan ke publik.

\item \textbf{Platform Target:}
Verifin adalah aplikasi web yang dioptimalkan untuk desktop browser (Chrome, Firefox, Safari).
Tidak ada aplikasi mobile native dalam cakupan pengembangan ini, meskipun antarmuka bersifat
responsif untuk mobile.

\item \textbf{Keterbatasan Model NLP:}
Model NLP berbasis IndoBERT yang digunakan adalah fine-tuned pada dataset lowongan kerja
sintetis dan publik yang tersedia. Akurasi model pada lowongan dengan pola penipuan yang sangat
baru atau tidak umum mungkin lebih rendah.

\item \textbf{Ketergantungan API Eksternal:}
Pipeline analisis bergantung pada layanan eksternal (OpenAgentic/OpenAI untuk LLM, Supabase
untuk database). Gangguan pada layanan ini akan mempengaruhi ketersediaan sistem.

\item \textbf{Bukan Layanan Hukum:}
Verdict yang dihasilkan Verifin adalah \textbf{penilaian berbasis risiko}, bukan keputusan
hukum. Sistem tidak dapat memastikan secara definitif bahwa sebuah lowongan adalah penipuan.
Pengguna tetap dianjurkan untuk melakukan verifikasi tambahan pada lowongan dengan skor
perbatasan.

\item \textbf{Privasi dan Anonimisasi:}
Laporan komunitas disimpan dengan IP yang di-hash, bukan IP asli. Data nomor telepon dan
informasi pribadi dari teks lowongan yang dianalisis tidak diekspos ke antarmuka publik.

\end{enumerate}


\section*{D.\quad Metodologi Pengembangan Perangkat Lunak}
\addcontentsline{toc}{section}{D. Metodologi Pengembangan Perangkat Lunak}

Verifin dikembangkan menggunakan metodologi \textbf{Agile Scrum} yang diadaptasi untuk tim kecil
(3 orang) dalam konteks kompetisi dengan \textit{time-box} ketat. Pengembangan dibagi menjadi
empat sprint dengan durasi dua minggu per sprint.

\subsection*{D.1 Struktur Sprint}

\begin{longtable}{@{}p{2.5cm}p{3.5cm}p{8cm}@{}}
\toprule
\textbf{Sprint} & \textbf{Periode} & \textbf{Target Deliverable} \\
\midrule
\endfirsthead
\toprule
\textbf{Sprint} & \textbf{Periode} & \textbf{Target Deliverable} \\
\midrule
\endhead
Sprint 0 & Minggu 1--2 & Riset, desain arsitektur, setup repo, konfigurasi Supabase, scaffolding Next.js dan FastAPI \\
\midrule
Sprint 1 & Minggu 3--4 & Implementasi pipeline NER + NLP Classifier, integrasi IndoBERT, endpoint FastAPI dasar, UI input form \\
\midrule
Sprint 2 & Minggu 5--6 & Implementasi modul OSINT (domain, phone, company), LLM synthesis via OpenAgentic, Fraud Network Graph, XAI explainer \\
\midrule
Sprint 3 & Minggu 7--8 & Trust Score engine, Community Monitoring, riwayat analisis, API publik, UI result dashboard, visualisasi graf \\
\midrule
Sprint 4 & Minggu 9--10 & Pengujian integrasi, optimasi performa, deployment Vercel + Railway, dokumentasi, persiapan demonstrasi \\
\bottomrule
\end{longtable}

\subsection*{D.2 Pembagian Peran Tim}

\begin{longtable}{@{}p{5cm}p{4cm}p{5.5cm}@{}}
\toprule
\textbf{Peran} & \textbf{Anggota} & \textbf{Tanggung Jawab Utama} \\
\midrule
\endfirsthead
\toprule
\textbf{Peran} & \textbf{Anggota} & \textbf{Tanggung Jawab Utama} \\
\midrule
\endhead
Tech Lead \& AI Engineer & Hafidz Rizqullah P. & Arsitektur sistem, pipeline NER+NLP, Trust Score engine, XAI, koordinasi tim \\
\midrule
Backend \& OSINT Engineer & Akmal Manggala P. & FastAPI backend, modul OSINT, Fraud Network Graph, database schema, API publik \\
\midrule
Frontend \& Integration & Matthew Hayunaji P. & Next.js frontend, UI/UX, visualisasi graf, integrasi LLM, community monitoring \\
\bottomrule
\end{longtable}

\subsection*{D.3 Prinsip Pengembangan}

\begin{enumerate}[leftmargin=*, label=\textbf{\arabic*.}]
\item \textbf{API-First:} Backend dirancang sebagai API murni sehingga frontend dan pihak ketiga
      dapat mengintegrasikan layanan dengan mudah.
\item \textbf{Test-Driven:} Setiap modul pipeline dilengkapi unit test sebelum integrasi.
\item \textbf{Security by Design:} Validasi input, rate limiting, dan manajemen secret diterapkan
      sejak Sprint 0.
\item \textbf{Explainability First:} Setiap keputusan Trust Score harus dapat dijelaskan dengan
      kontribusi fitur yang terukur.
\item \textbf{Progressive Enhancement:} Sistem fungsional minimal tersedia di akhir setiap sprint;
      tidak ada sprint yang berakhir tanpa \textit{working software}.
\end{enumerate}

\subsection*{D.4 Tools dan Platform Pengembangan}

\begin{longtable}{@{}p{4cm}p{4cm}p{6.5cm}@{}}
\toprule
\textbf{Kategori} & \textbf{Tool} & \textbf{Penggunaan} \\
\midrule
\endfirsthead
\toprule
\textbf{Kategori} & \textbf{Tool} & \textbf{Penggunaan} \\
\midrule
\endhead
Version Control & Git + GitHub & Branching model: main/develop/feature/* \\
Project Management & GitHub Projects & Kanban board, issue tracking, sprint planning \\
CI/CD & GitHub Actions & Lint, test, dan deploy otomatis ke staging \\
Komunikasi & Discord & Daily standup async, code review notifikasi \\
Dokumentasi & Notion & ADR (Architecture Decision Records), API docs draft \\
Testing & pytest + Jest & Unit test backend (pytest) dan frontend (Jest/Vitest) \\
\bottomrule
\end{longtable}


\section*{E.\quad Analisis Kebutuhan dan Desain Solusi Perangkat Lunak}
\addcontentsline{toc}{section}{E. Analisis Kebutuhan dan Desain Solusi Perangkat Lunak}

\subsection*{E.1 Analisis Kebutuhan Fungsional}

\begin{longtable}{@{}p{2cm}p{4cm}p{7.5cm}p{1.5cm}@{}}
\toprule
\textbf{ID} & \textbf{Nama} & \textbf{Deskripsi} & \textbf{Prioritas} \\
\midrule
\endfirsthead
\toprule
\textbf{ID} & \textbf{Nama} & \textbf{Deskripsi} & \textbf{Prioritas} \\
\midrule
\endhead
FR-01 & Input Teks Lowongan & Pengguna dapat menempelkan teks lowongan kerja (min. 50 karakter) ke dalam form input dan mengirimkan untuk dianalisis & Tinggi \\
\midrule
FR-02 & Input URL Lowongan & Pengguna dapat memasukkan URL halaman lowongan kerja; sistem mengekstrak teks dari halaman tersebut secara otomatis & Tinggi \\
\midrule
FR-03 & Ekstraksi Entitas NER & Sistem mengekstrak entitas: nama perusahaan, nomor telepon, email, URL/domain, nama kontak, posisi, gaji, lokasi, dan jenis pekerjaan & Tinggi \\
\midrule
FR-04 & Klasifikasi NLP & Sistem mengklasifikasikan teks lowongan ke dalam kategori LEGIT/SUSPICIOUS/SCAM menggunakan model fine-tuned IndoBERT & Tinggi \\
\midrule
FR-05 & Investigasi OSINT Domain & Sistem melakukan investigasi terhadap domain/URL yang ditemukan: Whois, usia domain, reputasi VirusTotal & Tinggi \\
\midrule
FR-06 & Investigasi OSINT Telepon & Sistem memvalidasi format dan prefix nomor telepon Indonesia; mengecek kemunculan nomor di database laporan komunitas & Tinggi \\
\midrule
FR-07 & Investigasi OSINT Perusahaan & Sistem mencari nama perusahaan di sumber publik dan mengidentifikasi inkonsistensi dengan klaim dalam lowongan & Tinggi \\
\midrule
FR-08 & Sintesis Naratif LLM & Sistem menghasilkan ringkasan narasi berbahasa Indonesia menggunakan LLM (GPT-4o-mini) berdasarkan temuan OSINT & Tinggi \\
\midrule
FR-09 & Trust Score dan Verdict & Sistem menghasilkan Trust Score integer 0--100 beserta verdict tiga tingkat: AMAN (61--100), WASPADA (31--60), BAHAYA (0--30) & Tinggi \\
\midrule
FR-10 & Penjelasan XAI & Sistem menghasilkan penjelasan kontribusi fitur terhadap Trust Score menggunakan custom SHAP-inspired additive explainer & Tinggi \\
\midrule
FR-11 & Visualisasi Graf Jaringan & Sistem menampilkan graf interaktif koneksi entitas lowongan dengan entitas lain di database & Sedang \\
\midrule
FR-12 & Community Monitoring & Pengguna dapat melaporkan lowongan mencurigakan, memberi upvote/downvote laporan, dan melihat daftar laporan & Sedang \\
\midrule
FR-13 & Riwayat Analisis & Sistem menyimpan riwayat analisis pengguna berdasarkan sesi/akun & Sedang \\
\midrule
FR-14 & API Publik & Sistem menyediakan REST API terdokumentasi via FastAPI auto-docs untuk integrasi pihak ketiga & Sedang \\
\bottomrule
\caption{Daftar Kebutuhan Fungsional Verifin}
\end{longtable}

\subsection*{E.2 Analisis Kebutuhan Non-Fungsional}

\begin{longtable}{@{}p{2cm}p{3.5cm}p{9cm}@{}}
\toprule
\textbf{ID} & \textbf{Nama} & \textbf{Deskripsi} \\
\midrule
\endfirsthead
\toprule
\textbf{ID} & \textbf{Nama} & \textbf{Deskripsi} \\
\midrule
\endhead
NFR-01 & Performa & Pipeline analisis lengkap (NER $\to$ NLP $\to$ OSINT $\to$ LLM $\to$ XAI) selesai maksimal 30 detik untuk 90\% permintaan. Modul OSINT dijalankan paralel. \\
\midrule
NFR-02 & Ketersediaan & Uptime minimal 95\% selama periode demonstrasi dan evaluasi. \\
\midrule
NFR-03 & Keamanan & Input divalidasi dan disanitasi; API key disimpan sebagai environment variable; rate limiting pada endpoint analisis. \\
\midrule
NFR-04 & Skalabilitas & Arsitektur FastAPI memungkinkan horizontal scaling; PostgreSQL di Supabase mendukung koneksi pooling. \\
\midrule
NFR-05 & Kemudahan Penggunaan & Dapat digunakan pencari kerja awam tanpa pelatihan teknis; proses analisis cukup paste teks dan klik tombol. \\
\midrule
NFR-06 & Transparansi & Setiap verdict disertai penjelasan berbasis fitur yang dapat dipahami dan diverifikasi. \\
\bottomrule
\caption{Daftar Kebutuhan Non-Fungsional Verifin}
\end{longtable}


\subsection*{E.3 Desain Solusi: Arsitektur Sistem}

Verifin mengimplementasikan arsitektur berlapis (\textit{layered architecture}) yang memisahkan
tanggung jawab setiap komponen secara jelas. Berikut adalah diagram arsitektur sistem secara keseluruhan:

\vspace{0.5cm}

% ── DIAGRAM PIPELINE 5-LAYER ──────────────────────────────────────────────
\begin{figure}[H]
\centering
\begin{tikzpicture}[
  node distance=0.6cm,
  box/.style={rectangle, rounded corners=4pt, minimum width=2.2cm, minimum height=1.1cm,
              text centered, font=\small\bfseries, text width=2.1cm},
  arrow/.style={-{Stealth[length=6pt]}, line width=1.5pt, color=verfinblue},
]
% Nodes
\node[box, fill=verfinblue!90, text=white] (ner)
      {Layer 1\\NER\\Ekstraksi Entitas};
\node[box, fill=verfinblue!70, text=white, right=of ner] (nlp)
      {Layer 2\\NLP\\Klasifikasi};
\node[box, fill=verfinblue!50, text=white, right=of nlp] (osint)
      {Layer 3\\OSINT\\Investigasi};
\node[box, fill=verfinblue!30, text=verifindark, right=of osint] (llm)
      {Layer 4\\LLM\\Sintesis};
\node[box, fill=verfingreen!70, text=white, right=of llm] (xai)
      {Layer 5\\XAI\\Penjelasan};

% Arrows
\draw[arrow] (ner) -- (nlp);
\draw[arrow] (nlp) -- (osint);
\draw[arrow] (osint) -- (llm);
\draw[arrow] (llm) -- (xai);

% Labels below
\node[font=\tiny, below=0.15cm of ner, text=gray] {IndoBERT NER};
\node[font=\tiny, below=0.15cm of nlp, text=gray] {IndoBERT Cls};
\node[font=\tiny, below=0.15cm of osint, text=gray] {Whois/Phone/Co.};
\node[font=\tiny, below=0.15cm of llm, text=gray] {GPT-4o-mini};
\node[font=\tiny, below=0.15cm of xai, text=gray] {SHAP-inspired};

% Input arrow
\node[left=0.5cm of ner, font=\small\itshape, text=verifindark, align=center, text width=1.5cm] (input) {Input\\Lowongan};
\draw[arrow] (input) -- (ner);

% Output arrow
\node[right=0.5cm of xai, font=\small\bfseries, text=verfingreen, align=center, text width=1.5cm] (output) {Trust\\Score};
\draw[arrow] (xai) -- (output);
\end{tikzpicture}
\caption{Pipeline analisis 5-layer Verifin: dari input teks hingga Trust Score dan XAI}
\end{figure}

\vspace{0.3cm}

% ── DIAGRAM JOB TRUST INFRASTRUCTURE ─────────────────────────────────────
\begin{figure}[H]
\centering
\begin{tikzpicture}[
  every node/.style={font=\small},
]
% Outer box — Job Trust Infrastructure
\node[rectangle, draw=verfinblue, line width=2pt, rounded corners=8pt,
      minimum width=12cm, minimum height=6.5cm,
      label={[font=\bfseries\normalsize,color=verfinblue]above:Job Trust Infrastructure}] (outer) {};

% Inner top-left: OSINT Engine
\node[rectangle, draw=verfinblue!60, fill=verfinblue!10, rounded corners=4pt,
      minimum width=4.5cm, minimum height=2.2cm, text centered,
      at={($(outer.north west)+(2.7cm,-1.7cm)$)}] (osint_box) {};
\node[font=\bfseries\small, color=verfinblue] at ($(osint_box.north)+(0,-0.3cm)$) {OSINT Engine};
\node[font=\tiny, align=center] at ($(osint_box.center)+(0,-0.1cm)$)
      {Domain Whois \\ Phone Validator \\ Company Lookup \\ VirusTotal API};

% Inner top-right: Fraud Network Graph
\node[rectangle, draw=verfinred!70, fill=verfinred!10, rounded corners=4pt,
      minimum width=4.5cm, minimum height=2.2cm, text centered,
      at={($(outer.north east)+(-2.7cm,-1.7cm)$)}] (fraud_box) {};
\node[font=\bfseries\small, color=verfinred] at ($(fraud_box.north)+(0,-0.3cm)$) {Fraud Network Graph};
\node[font=\tiny, align=center] at ($(fraud_box.center)+(0,-0.1cm)$)
      {PostgreSQL Graph DB \\ SHA-256 Fingerprint \\ Entity Deduplication \\ Community Reports};

% Inner bottom-left: Trust Scoring
\node[rectangle, draw=verfinyellow!80, fill=verfinyellow!10, rounded corners=4pt,
      minimum width=4.5cm, minimum height=2.2cm, text centered,
      at={($(outer.south west)+(2.7cm,1.7cm)$)}] (trust_box) {};
\node[font=\bfseries\small, color=verfinyellow!80!black] at ($(trust_box.north)+(0,-0.3cm)$) {Trust Scoring Engine};
\node[font=\tiny, align=center] at ($(trust_box.center)+(0,-0.1cm)$)
      {Weighted Aggregation \\ NLP + OSINT + Community \\ Verdict Assignment \\ 0--100 Integer Score};

% Inner bottom-right: XAI Explainer
\node[rectangle, draw=verfingreen!70, fill=verfingreen!10, rounded corners=4pt,
      minimum width=4.5cm, minimum height=2.2cm, text centered,
      at={($(outer.south east)+(-2.7cm,1.7cm)$)}] (xai_box) {};
\node[font=\bfseries\small, color=verfingreen!70!black] at ($(xai_box.north)+(0,-0.3cm)$) {XAI Explainer};
\node[font=\tiny, align=center] at ($(xai_box.center)+(0,-0.1cm)$)
      {SHAP-inspired Additive \\ Feature Contribution \\ Human-readable Output \\ LLM Narrative Synthesis};

% Arrows connecting components
\draw[-{Stealth}, color=verfinblue!60, line width=1pt] (osint_box) -- (trust_box);
\draw[-{Stealth}, color=verfinred!60, line width=1pt] (fraud_box) -- (trust_box);
\draw[-{Stealth}, color=verfinyellow!80, line width=1pt] (trust_box) -- (xai_box);
\draw[-{Stealth}, color=verfinblue!60, dashed, line width=1pt] (osint_box) -- (fraud_box);
\end{tikzpicture}
\caption{Job Trust Infrastructure: komponen inti sistem Verifin}
\end{figure}

\vspace{0.3cm}

% ── DIAGRAM VERDICT TIERS ─────────────────────────────────────────────────
\begin{figure}[H]
\centering
\begin{tikzpicture}[every node/.style={font=\normalsize}]
% BAHAYA box
\node[rectangle, draw=verfinred, fill=verfinred!15, line width=2pt, rounded corners=6pt,
      minimum width=3.5cm, minimum height=2.5cm, text centered] (bahaya) at (0,0) {};
\node[font=\Large\bfseries, color=verfinred] at ($(bahaya.center)+(0,0.4cm)$) {BAHAYA};
\node[font=\small, color=verfinred!80!black] at ($(bahaya.center)+(0,-0.1cm)$) {Skor 0--30};
\node[font=\tiny, color=gray, text width=3cm, align=center] at ($(bahaya.center)+(0,-0.6cm)$)
      {Indikator kuat penipuan};

% WASPADA box
\node[rectangle, draw=verfinyellow!80!black, fill=verfinyellow!20, line width=2pt, rounded corners=6pt,
      minimum width=3.5cm, minimum height=2.5cm, text centered] (waspada) at (4.5,0) {};
\node[font=\Large\bfseries, color=verfinyellow!70!black] at ($(waspada.center)+(0,0.4cm)$) {WASPADA};
\node[font=\small, color=verfinyellow!70!black] at ($(waspada.center)+(0,-0.1cm)$) {Skor 31--60};
\node[font=\tiny, color=gray, text width=3cm, align=center] at ($(waspada.center)+(0,-0.6cm)$)
      {Sinyal meragukan, bukti tidak konklusif};

% AMAN box
\node[rectangle, draw=verfingreen, fill=verfingreen!15, line width=2pt, rounded corners=6pt,
      minimum width=3.5cm, minimum height=2.5cm, text centered] (aman) at (9,0) {};
\node[font=\Large\bfseries, color=verfingreen!80!black] at ($(aman.center)+(0,0.4cm)$) {AMAN};
\node[font=\small, color=verfingreen!80!black] at ($(aman.center)+(0,-0.1cm)$) {Skor 61--100};
\node[font=\tiny, color=gray, text width=3cm, align=center] at ($(aman.center)+(0,-0.6cm)$)
      {Indikator kepercayaan kuat dari OSINT};

% Scale bar below
\shade[left color=verfinred!60, right color=verfingreen!60]
      (-1.7cm,-1.8cm) rectangle (10.7cm,-1.4cm);
\node[font=\tiny, color=white] at (0,-1.6cm) {0};
\node[font=\tiny, color=white] at (4.5cm,-1.6cm) {50};
\node[font=\tiny, color=white] at (9cm,-1.6cm) {100};
\end{tikzpicture}
\caption{Tiga tingkat verdict Verifin berdasarkan rentang Trust Score}
\end{figure}

\vspace{0.3cm}

% ── TRUST SCORE TABLE ─────────────────────────────────────────────────────
\begin{table}[H]
\centering
\caption{Tabel Trust Score dan Interpretasi Verdict}
\begin{tabular}{@{}p{2.5cm}p{2.5cm}p{8cm}@{}}
\toprule
\textbf{Rentang Skor} & \textbf{Verdict} & \textbf{Interpretasi} \\
\midrule
\rowcolor{verfinred!20}
0--30 & \textbf{\color{verfinred}BAHAYA} & Terdapat indikator kuat penipuan dari verifikasi OSINT: domain baru/tidak valid, nomor tidak terdaftar, nama perusahaan tidak ditemukan, pola teks sesuai ciri penipuan. Pengguna sangat disarankan tidak merespons. \\
\midrule
\rowcolor{verfinyellow!20}
31--60 & \textbf{\color{verfinyellow!70!black}WASPADA} & Terdapat beberapa sinyal meragukan namun bukti tidak konklusif. Pengguna disarankan melakukan verifikasi tambahan sebelum merespons. \\
\midrule
\rowcolor{verfingreen!15}
61--100 & \textbf{\color{verfingreen!70!black}AMAN} & Lowongan memiliki indikator kepercayaan yang kuat dari verifikasi OSINT: domain valid dan berumur, perusahaan terverifikasi, tidak ada laporan penipuan. \\
\bottomrule
\end{tabular}
\end{table}


\subsection*{E.4 Desain Solusi: Tech Stack}

\begin{longtable}{@{}p{3.5cm}p{4cm}p{7cm}@{}}
\toprule
\textbf{Komponen} & \textbf{Teknologi} & \textbf{Keterangan} \\
\midrule
\endfirsthead
\toprule
\textbf{Komponen} & \textbf{Teknologi} & \textbf{Keterangan} \\
\midrule
\endhead
Frontend Framework & Next.js 14 (App Router) & React SSR/SSG, routing berbasis file, TypeScript \\
\midrule
UI Styling & Tailwind CSS 3 & Utility-first CSS, kustomisasi penuh \\
\midrule
Animasi UI & motion/react & Animasi deklaratif, transisi halaman \\
\midrule
Ikon & Phosphor Icons & Library ikon konsisten dan ringan \\
\midrule
Visualisasi Graf & React Flow + D3.js & Graf interaktif Fraud Network \\
\midrule
Backend Framework & FastAPI (Python) & REST API async, auto OpenAPI docs \\
\midrule
NER Model & IndoBERT (fine-tuned) & Named Entity Recognition Bahasa Indonesia \\
\midrule
NLP Classifier & IndoBERT (fine-tuned) & Klasifikasi LEGIT/SUSPICIOUS/SCAM \\
\midrule
LLM Orchestration & OpenAgentic + GPT-4o-mini & Sintesis narasi berbasis bukti OSINT \\
\midrule
XAI Engine & Custom SHAP-inspired & Additive feature contribution explainer \\
\midrule
OSINT: Domain & python-whois, RDAP & Whois lookup, usia domain \\
\midrule
OSINT: Reputasi & VirusTotal API & Cek reputasi domain/IP \\
\midrule
OSINT: Perusahaan & Google Custom Search API & Pencarian nama perusahaan publik \\
\midrule
Database & PostgreSQL (Supabase) & Data lowongan, laporan, fingerprint \\
\midrule
ORM & SQLAlchemy + Alembic & Query ORM, migrasi schema \\
\midrule
Auth & Supabase Auth & JWT-based, opsional untuk komunitas \\
\midrule
Deployment Frontend & Vercel & CDN global, preview deployments \\
\midrule
Deployment Backend & Railway & Container Python, auto-scaling \\
\midrule
CI/CD & GitHub Actions & Lint, test, deploy otomatis \\
\midrule
Testing Backend & pytest & Unit dan integration test \\
\midrule
Testing Frontend & Jest + Vitest & Komponen dan hook test \\
\bottomrule
\caption{Tech Stack Verifin}
\end{longtable}


\section*{F.\quad Implementasi Perangkat Lunak}
\addcontentsline{toc}{section}{F. Implementasi Perangkat Lunak}

\subsection*{F.1 Implementasi Pipeline Analisis}

\subsubsection*{Layer 1: Named Entity Recognition (NER)}

Sistem menggunakan model IndoBERT yang di-\textit{fine-tune} pada dataset lowongan kerja Bahasa
Indonesia untuk mengekstrak entitas-entitas kunci berikut:

\begin{itemize}[leftmargin=*]
\item \textbf{ORG} --- Nama perusahaan atau organisasi
\item \textbf{PHONE} --- Nomor telepon/WhatsApp
\item \textbf{EMAIL} --- Alamat surel
\item \textbf{URL} --- Alamat website atau domain
\item \textbf{PERSON} --- Nama kontak person (HRD, rekruter)
\item \textbf{SALARY} --- Informasi gaji (nominal, range)
\item \textbf{LOC} --- Lokasi kerja
\item \textbf{POSITION} --- Posisi/jabatan yang ditawarkan
\item \textbf{WORKTYPE} --- Jenis kerja (WFH, onsite, hybrid)
\end{itemize}

Model NER diimplementasikan menggunakan pustaka \texttt{transformers} (HuggingFace) dengan
tokenizer IndoBERT. Output berupa JSON terstruktur yang menjadi input untuk semua layer berikutnya.

\subsubsection*{Layer 2: NLP Classifier}

Model klasifikasi berbasis IndoBERT mengklasifikasikan teks lowongan ke dalam tiga kategori:
\texttt{LEGIT}, \texttt{SUSPICIOUS}, atau \texttt{SCAM}. Model menggunakan representasi
\texttt{[CLS]} token dari IndoBERT yang dipetakan ke tiga kelas melalui \textit{classification head}
linear. Skor probabilitas dari model ini menjadi salah satu fitur utama dalam kalkulasi Trust Score.

Pola-pola yang menjadi sinyal mencurigakan dalam model antara lain:
\begin{itemize}[leftmargin=*]
\item Gaji tidak realistis untuk syarat minimal
\item Permintaan data pribadi (KTP, foto, rekening) di awal proses
\item Tidak ada informasi kontak perusahaan yang jelas
\item Janji kerja langsung tanpa proses seleksi
\item Desakan waktu yang tidak wajar (\textit{limited slot}, \textit{closing soon})
\item Tidak ada deskripsi pekerjaan yang spesifik
\end{itemize}

\subsubsection*{Layer 3: OSINT Engine}

Modul OSINT dijalankan secara paralel menggunakan \texttt{asyncio} untuk meminimalkan latensi.
Terdapat tiga sub-modul utama:

\textbf{Domain/URL Investigator:}
\begin{itemize}[leftmargin=*]
\item Whois/RDAP lookup untuk mendapatkan tanggal registrasi, registrar, dan informasi pemilik
\item Domain berumur $< 30$ hari mendapat penalti skor signifikan
\item Validasi SSL/HTTPS dan kesesuaian nama domain dengan klaim perusahaan
\item Pengecekan reputasi via VirusTotal API (domain dan IP)
\end{itemize}

\textbf{Phone Number Validator:}
\begin{itemize}[leftmargin=*]
\item Validasi format nomor telepon Indonesia (08xx, +62xx)
\item Identifikasi operator berdasarkan prefix (Telkomsel, Indosat, XL, dll.)
\item Pengecekan apakah nomor pernah dilaporkan di database komunitas Verifin
\item Flag khusus untuk nomor dengan banyak laporan penipuan
\end{itemize}

\textbf{Company Lookup:}
\begin{itemize}[leftmargin=*]
\item Pencarian nama perusahaan via Google Custom Search API dengan kueri terstruktur
\item Verifikasi keberadaan website resmi dan kesesuaian domain
\item Deteksi inkonsistensi antara klaim jabatan/lokasi dengan temuan publik
\end{itemize}

\subsubsection*{Layer 4: LLM Narrative Synthesis}

Sistem mengirimkan bundle hasil OSINT ke GPT-4o-mini melalui OpenAgentic dengan prompt terstruktur
yang memerintahkan model untuk:
\begin{enumerate}[leftmargin=*]
\item Merangkum temuan OSINT dalam narasi berbahasa Indonesia yang mudah dipahami
\item Menyebutkan fakta-fakta spesifik yang mendukung atau melemahkan kepercayaan
\item Memberikan rekomendasi tindak lanjut yang actionable
\item Tidak membuat klaim yang tidak didukung oleh data OSINT yang tersedia
\end{enumerate}

Penggunaan LLM dibatasi hanya untuk sintesis narasi; keputusan Trust Score tidak bergantung
pada output LLM untuk menghindari halusinasi yang mempengaruhi verdict.

\subsubsection*{Layer 5: XAI Explainer}

Custom SHAP-inspired additive feature explainer menghitung kontribusi setiap fitur terhadap
Trust Score final. Fitur-fitur yang dipertimbangkan antara lain:

\begin{longtable}{@{}p{5cm}p{3cm}p{6.5cm}@{}}
\toprule
\textbf{Fitur} & \textbf{Bobot Dasar} & \textbf{Deskripsi} \\
\midrule
\endfirsthead
\toprule
\textbf{Fitur} & \textbf{Bobot Dasar} & \textbf{Deskripsi} \\
\midrule
\endhead
NLP Classifier Score & 25\% & Probabilitas LEGIT dari model IndoBERT \\
Domain Age & 15\% & Usia domain: makin tua makin tinggi skor \\
Domain Reputation & 15\% & Skor VirusTotal: clean = bonus, malicious = penalti \\
Phone Reported & 10\% & Nomor pernah dilaporkan = penalti besar \\
Company Found & 15\% & Perusahaan terverifikasi di sumber publik \\
Fraud Network Match & 20\% & Entitas cocok dengan fingerprint penipuan sebelumnya \\
\bottomrule
\caption{Fitur dan bobot dasar XAI Explainer Verifin}
\end{longtable}

Output XAI berupa daftar fitur beserta nilai kontribusinya (positif/negatif) dan penjelasan
dalam bahasa Indonesia yang dapat dipahami pengguna awam.

\subsection*{F.2 Struktur Kode}

Repositori Verifin diorganisasikan sebagai monorepo dengan struktur berikut:

\begin{lstlisting}[language=bash, caption={Struktur direktori repositori Verifin}]
verifin/
|- backend/
|   |- app/
|   |   |- main.py              # FastAPI entry point
|   |   |- routers/
|   |   |   |- analyze.py       # POST /analyze endpoint
|   |   |   |- community.py     # GET/POST /community endpoints
|   |   |   |- history.py       # GET /history endpoint
|   |   |- services/
|   |   |   |- ner_service.py   # IndoBERT NER pipeline
|   |   |   |- nlp_service.py   # IndoBERT classifier
|   |   |   |- osint/
|   |   |   |   |- domain.py    # Domain/Whois investigator
|   |   |   |   |- phone.py     # Phone validator
|   |   |   |   |- company.py   # Company lookup
|   |   |   |- llm_service.py   # OpenAgentic LLM synthesis
|   |   |   |- xai_service.py   # SHAP-inspired explainer
|   |   |   |- trust_score.py   # Score aggregation engine
|   |   |   |- graph_service.py # Fraud Network Graph queries
|   |   |- models/
|   |   |   |- job.py           # SQLAlchemy Job model
|   |   |   |- report.py        # Community report model
|   |   |   |- fingerprint.py   # Fraud fingerprint model
|   |   |- database.py          # Supabase/PostgreSQL connection
|   |- tests/                   # pytest unit & integration tests
|   |- requirements.txt
|- frontend/
|   |- src/
|   |   |- app/
|   |   |   |- page.tsx         # Halaman utama / landing
|   |   |   |- analyze/page.tsx # Halaman analisis
|   |   |   |- result/[id]/page.tsx  # Halaman hasil analisis
|   |   |   |- community/page.tsx    # Halaman komunitas
|   |   |- components/
|   |   |   |- AnalysisForm.tsx       # Form input teks/URL
|   |   |   |- VerdictBadge.tsx       # AMAN/WASPADA/BAHAYA badge
|   |   |   |- TrustScoreGauge.tsx    # Visualisasi skor 0-100
|   |   |   |- OsintBreakdown.tsx     # Detail hasil per modul OSINT
|   |   |   |- XaiExplainer.tsx       # Breakdown kontribusi fitur
|   |   |   |- FraudNetworkGraph.tsx  # Visualisasi graf interaktif
|   |   |   |- CommunityFeed.tsx      # Feed laporan komunitas
|   |   |- lib/
|   |   |   |- api.ts           # API client functions
|   |   |   |- types.ts         # TypeScript type definitions
|   |- package.json
|   |- tailwind.config.ts
|- README.md
\end{lstlisting}

\subsection*{F.3 Skema Database PostgreSQL}

\textbf{Tabel \texttt{jobs}} --- Menyimpan data lowongan dan hasil analisis:

\begin{lstlisting}[language=SQL, caption={DDL tabel jobs}]
CREATE TABLE jobs (
    id            UUID PRIMARY KEY DEFAULT gen_random_uuid(),
    input_text    TEXT NOT NULL,
    input_url     TEXT,
    created_at    TIMESTAMPTZ DEFAULT NOW(),
    status        VARCHAR(20) DEFAULT 'pending', -- pending|processing|done|error
    trust_score   INTEGER,                        -- 0-100
    verdict       VARCHAR(10),                    -- AMAN|WASPADA|BAHAYA
    fingerprint   VARCHAR(64),                    -- SHA-256 hash entitas kunci
    entities      JSONB,                          -- Hasil ekstraksi NER
    nlp_score     FLOAT,                          -- Skor klasifikasi NLP
    osint_bundle  JSONB,                          -- Hasil seluruh modul OSINT
    llm_summary   TEXT,                           -- Narasi LLM berbasis bukti
    xai_factors   JSONB                           -- Kontribusi fitur XAI
);
\end{lstlisting}

\textbf{Tabel \texttt{community\_reports}} --- Menyimpan laporan pengguna:

\begin{lstlisting}[language=SQL, caption={DDL tabel community\_reports}]
CREATE TABLE community_reports (
    id            UUID PRIMARY KEY DEFAULT gen_random_uuid(),
    job_id        UUID REFERENCES jobs(id),
    reporter_ip   VARCHAR(45),        -- Hashed IP
    report_text   TEXT,
    upvotes       INTEGER DEFAULT 0,
    downvotes     INTEGER DEFAULT 0,
    created_at    TIMESTAMPTZ DEFAULT NOW(),
    verified      BOOLEAN DEFAULT FALSE
);
\end{lstlisting}

\textbf{Tabel \texttt{fraud\_fingerprints}} --- Indeks deduplikasi:

\begin{lstlisting}[language=SQL, caption={DDL tabel fraud\_fingerprints}]
CREATE TABLE fraud_fingerprints (
    fingerprint   VARCHAR(64) PRIMARY KEY,
    first_seen    TIMESTAMPTZ DEFAULT NOW(),
    last_seen     TIMESTAMPTZ DEFAULT NOW(),
    occurrence    INTEGER DEFAULT 1,
    avg_score     FLOAT,
    entity_hash   JSONB  -- Komponen hash per entitas
);
\end{lstlisting}

\subsection*{F.4 Contoh Alur Analisis End-to-End}

Berikut adalah ilustrasi alur kerja sistem ketika pengguna mengirimkan teks lowongan:

\textbf{Input dari pengguna:}

\begin{lstlisting}[caption={Contoh teks lowongan yang dianalisis}]
Dibutuhkan Admin Online WFH
Gaji 8-15 juta/bulan
Syarat: min SMA, punya HP Android
Hubungi: 0812-XXXX-XXXX (Budi - HRD PT Sukses Mandiri)
Langsung kerja, tidak perlu pengalaman
Kirim foto dan KTP sekarang
\end{lstlisting}

\textbf{Layer 1 --- NER Output:}

\begin{lstlisting}[language=Python, caption={Output JSON NER Layer 1}]
{
  "company": "PT Sukses Mandiri",
  "contact_phone": "0812XXXXXXXX",
  "contact_name": "Budi",
  "role": "HRD",
  "position": "Admin Online",
  "work_type": "WFH",
  "salary_min": 8000000,
  "salary_max": 15000000,
  "requirements": ["min SMA", "HP Android"],
  "red_flags_text": ["kirim foto", "KTP sekarang",
                     "langsung kerja", "tidak perlu pengalaman"]
}
\end{lstlisting}

\textbf{Layer 2 --- NLP Output:}

\begin{lstlisting}[language=Python, caption={Output JSON NLP Classifier Layer 2}]
{
  "label": "SCAM",
  "confidence": 0.87,
  "probabilities": {"LEGIT": 0.05, "SUSPICIOUS": 0.08, "SCAM": 0.87}
}
\end{lstlisting}

\textbf{Layer 3 --- OSINT Output:}

\begin{lstlisting}[language=Python, caption={Output JSON OSINT Layer 3}]
{
  "domain": null,
  "phone": {
    "number": "0812XXXXXXXX",
    "operator": "Telkomsel",
    "format_valid": true,
    "community_reports": 3,
    "last_report": "2026-06-15"
  },
  "company": {
    "name": "PT Sukses Mandiri",
    "found_public": false,
    "website_found": false,
    "inconsistencies": ["Nama terlalu generik, tidak ditemukan di sumber publik"]
  }
}
\end{lstlisting}

\textbf{Layer 4 --- LLM Narrative:}

\begin{quote}
\textit{``Lowongan ini menunjukkan beberapa pola yang sangat umum pada penipuan rekrutmen. Nomor
kontak yang diberikan (0812-XXXX-XXXX) telah dilaporkan 3 kali oleh komunitas Verifin sebagai
nomor penipuan. Nama perusahaan PT Sukses Mandiri tidak ditemukan dalam sumber publik manapun.
Permintaan KTP dan foto di awal proses tanpa verifikasi identitas rekruter adalah tanda bahaya
besar. Gaji 8--15 juta untuk posisi Admin Online WFH tanpa pengalaman tidak realistis untuk pasar
kerja Indonesia saat ini. Sangat disarankan untuk tidak merespons iklan ini.''}
\end{quote}

\textbf{Layer 5 --- XAI Output dan Trust Score Final:}

\begin{longtable}{@{}p{5cm}p{2.5cm}p{7cm}@{}}
\toprule
\textbf{Fitur} & \textbf{Kontribusi} & \textbf{Penjelasan} \\
\midrule
\endfirsthead
\toprule
\textbf{Fitur} & \textbf{Kontribusi} & \textbf{Penjelasan} \\
\midrule
\endhead
NLP Classifier & \textcolor{verfinred}{$-$22 poin} & Model mendeteksi pola teks SCAM (87\% confidence) \\
Phone Reports & \textcolor{verfinred}{$-$18 poin} & Nomor dilaporkan 3x sebagai penipuan \\
Company Not Found & \textcolor{verfinred}{$-$15 poin} & Perusahaan tidak ditemukan di sumber publik \\
No Domain & \textcolor{verfinred}{$-$8 poin} & Tidak ada domain/website yang bisa diverifikasi \\
Red Flag Text & \textcolor{verfinred}{$-$12 poin} & Permintaan KTP/foto awal + janji langsung kerja \\
Base Score & $+75$ & Skor dasar sebelum penilaian \\
\midrule
\textbf{Trust Score Final} & \textbf{\textcolor{verfinred}{0}} & \textbf{\textcolor{verfinred}{BAHAYA}} \\
\bottomrule
\caption{Breakdown XAI untuk contoh lowongan penipuan}
\end{longtable}

\subsection*{F.5 Biodata Anggota Tim}

\subsubsection*{Ketua Tim: Hafidz Rizqullah Prasetya}
\begin{tabular}{@{}p{5cm}p{9.5cm}@{}}
\textbf{NIM} & 24/535493/SV/24243 \\
\textbf{Fakultas/Sekolah} & Sekolah Vokasi \\
\textbf{Program Studi} & Teknologi Rekayasa Perangkat Lunak (D4) \\
\textbf{Tanggal Lahir} & 1 Januari 2006 \\
\textbf{Surel} & \texttt{hafidzrizqullahprasetya@mail.ugm.ac.id} \\
\textbf{No. HP/WA} & 081325081046 \\
\textbf{Alamat} & Gondangan Ringinsari RT 006/RW 050, Maguwoharjo, Depok, Sleman, D.I.\ Yogyakarta (55282) \\
\textbf{Pekerjaan Ortu} & Wiraswasta / Wiraswasta \\
\textbf{Kontak Ortu (Ibu)} & 081326810312 \\
\end{tabular}

\vspace{0.5cm}
\subsubsection*{Anggota Tim 1: Akmal Manggala Putra}
\begin{tabular}{@{}p{5cm}p{9.5cm}@{}}
\textbf{NIM} & 24/536182/SV/24402 \\
\textbf{Fakultas/Sekolah} & Sekolah Vokasi \\
\textbf{Program Studi} & Teknologi Rekayasa Perangkat Lunak (D4) \\
\textbf{Tanggal Lahir} & 28 Oktober 2005 \\
\textbf{Surel} & \texttt{akmalmanggala28@gmail.com} \\
\textbf{No. HP/WA} & 085870457382 \\
\textbf{Alamat} & Tambalan RT 05/RW 00, Srimartani, Piyungan, Bantul, D.I.\ Yogyakarta (55792) \\
\textbf{Pekerjaan Ortu} & Pensiunan PNS / Ibu Rumah Tangga \\
\textbf{Kontak Ortu} & 085643507464 \\
\end{tabular}

\vspace{0.5cm}
\subsubsection*{Anggota Tim 2: Matthew Hayunaji Priantara}
\begin{tabular}{@{}p{5cm}p{9.5cm}@{}}
\textbf{NIM} & 24/536179/SV/24400 \\
\textbf{Fakultas/Sekolah} & Sekolah Vokasi \\
\textbf{Program Studi} & Teknologi Rekayasa Perangkat Lunak (D4) \\
\textbf{Tanggal Lahir} & 7 Oktober 2005 \\
\textbf{Surel} & \texttt{matthewpriantara99@gmail.com} \\
\textbf{No. HP/WA} & 082223217875 \\
\textbf{Alamat} & Pundong 2 RT 02/RW 04, Tirtoadi, Mlati, Sleman, D.I.\ Yogyakarta (55284) \\
\textbf{Pekerjaan Ortu} & PNS / Ibu Rumah Tangga \\
\textbf{Kontak Ortu} & 082324286886 \\
\end{tabular}


\section*{Daftar Pustaka}
\addcontentsline{toc}{section}{Daftar Pustaka}

\begin{enumerate}[leftmargin=*, label={[\arabic*]}]

\item Badan Pusat Statistik (BPS), ``Keadaan Ketenagakerjaan Indonesia Februari 2025,''
  Berita Resmi Statistik No.\ 39/05/Th.\ XXVIII, Mei 2025.
  \url{https://www.bps.go.id}

\item Global Anti-Scam Alliance (GASA) \& GS Lab, ``Global State of Scams 2024,''
  GASA Annual Report, 2024.
  \url{https://www.gasa.org}

\item Komnas Perempuan, ``Laporan Tahunan Tindak Pidana Perdagangan Orang (TPPO) 2024,''
  Jakarta: Komnas Perempuan, 2024.

\item S.\ Shalini, R.\ Lokesh, and S.\ Priya, ``Fake Job Posting Detection Using Machine
  Learning,'' \textit{International Journal of Computer Applications}, vol.\ 183, no.\ 52,
  pp.\ 1--6, 2022.

\item T.\ Alwafi and R.\ Abdulrahman, ``Detection of Recruitment Fraud Using Semantic
  Approaches,'' \textit{Security and Communication Networks}, 2019,
  doi: \texttt{10.1155/2019/7523043}.

\item G.\ Carvallo and S.\ E.\ Benu, ``Tinjauan Yuridis Penanggulangan Cyber-Recruitment
  dalam Tindak Pidana Perdagangan Orang: Tinjauan Unsur dan Pembuktian Digital,''
  \textit{Majelis: Jurnal Hukum Indonesia}, vol.\ 3, no.\ 2, pp.\ 13--22, Mei 2026,
  doi: \texttt{10.62383/majelis.v3i2.1532}.

\item V.\ G.\ Varsha and P.\ A.\ Thomas, ``Explainable AI For Phishing Detection: Techniques,
  Challenges, and Experimental Validation,'' in \textit{Proc.\ IEEE Recent Advances in
  Intelligent Computational Systems (RAICS)}, Nov.\ 2025.

\item J.\ Pastor-Galindo, P.\ Nespoli, F.\ G.\ Marmol, and G.\ M.\ Perez, ``The Not Yet
  Exploited Goldmine of OSINT: Opportunities, Open Challenges and Future Trends,''
  \textit{IEEE Access}, vol.\ 8, pp.\ 10282--10304, 2020,
  doi: \texttt{10.1109/ACCESS.2020.2965257}.

\item B.\ Wilie et al., ``IndoNLU: Benchmark and Resources for Evaluating Indonesian Natural
  Language Understanding,'' in \textit{Proc.\ 1st Conf.\ of the Asia-Pacific Chapter of the
  ACL (AACL)}, Dec.\ 2020, pp.\ 843--857.

\end{enumerate}

\end{document}
